

\usepackage{fontspec}



% Redefine section and subsection commands to adjust heading sizes
\makeatletter
\renewcommand\section{\@startsection{section}{1}{\z@}%
    {-3.5ex \@plus -1ex \@minus -.2ex}%
    {2.3ex \@plus.2ex}%
    {\Huge\bfseries}}
\renewcommand\subsection{\@startsection{subsection}{2}{\z@}%
    {-3.25ex\@plus -1ex \@minus -.2ex}%
    {1.5ex \@plus .2ex}%
    {\huge\bfseries}}
\makeatother

\usepackage{xcolor}
\usepackage{graphicx}
\usepackage{tikz}  %# Required for positioning the image.
\AddToHook{shipout/foreground}{
  \begin{tikzpicture}[remember picture,overlay]
    \node[black,rotate=45,scale=12,opacity=0.05] at (current page.center) {Confidential}; 
  \end{tikzpicture}
}

\makeatletter  %# Customizing the title page.
\renewcommand{\maketitle}{
  \begin{titlepage}
    \centering
    \vspace*{\fill}
    % \begin{tikzpicture}[remember picture,overlay]
    %   \centering
    %   \node[anchor=north,inner sep=10pt] at (current page.north) {\includegraphics[width=\paperwidth,height=0.3\paperheight]{../fig/ins_wordcloud.png}}; % Adjust the path and dimensions of the image
    % \end{tikzpicture}
    \Huge\textbf{\@title} \\[1cm]  %# Title formatting.
    \Large\textbf{\@subtitle} \\[1cm]  %# Subtitle formatting.
    \textbf{\@date} \\[1cm]  %# Date formatting.
    \begin{tikzpicture}[remember picture,overlay]
      \node[anchor=south east,inner sep=0] at (current page.south east) {\includegraphics[width=15cm]{../fig/irem_logo.png}};  %# Logo positioning.
    \end{tikzpicture}
    \vspace*{\fill}
  \end{titlepage}
  \newpage
  \section*{Executive Summary}
  \Large\textbf{Overview}

\normalsize

The Institutional Research and Enrolment Management team at CNC conducted a survey on behalf of the Academic Success Centre to assess the effectiveness and reach of its testing and tutoring services. This survey targeted CNC students who were actively registered and enrolled in eligible courses (see Page 5 for cohort selection information). The survey's aim was to understand student awareness, usage, and satisfaction with the tutoring services provided by the Centre.

\Large\textbf{Survey Methodology}

\normalsize

The survey was conducted online using SurveyMonkey from February 13, 2024 to March 8, 2024, and it took participants approximately four minutes to complete. Participants were selected using a combination of random sampling and targeted sampling to ensure representation of underrepresented groups, such as Indigenous students. The invited cohort size was 1002 students, with a final sample size of 212 respondents, resulting in a response rate of 21.2\%. Participants had the chance to win gift cards as an incentive for their participation.

\Large\textbf{Key Findings}

\normalsize

    \textbf{Awareness and Usage:}
        A significant portion of the students were aware of the tutoring services offered by the Academic Success Centre, suggesting effective communication efforts by CNC.
        Despite awareness, 63\% of the respondents had not used the tutoring services. Reasons for not using the services included scheduling conflicts, lack of familiarity with how to access services, and a general sense of unease.

   \textbf{Communication Channels:}
        Students preferred learning about tutoring services through digital communications (emails, Moodle, and the CNC website). However, they most often reported first learning about tutoring services through in-person methods (e.g., in-class announcements, faculty/staff conversations, word of mouth).

    \textbf{Barriers to Access:}
        The most common barrier was scheduling conflicts. Students indicated they did not attempt to use the services due to lack of time or discomfort with the process.
        There is a clear need for more flexible and accessible tutoring options to accommodate students’ varied schedules.

    \textbf{Satisfaction Levels:}
        Overall satisfaction with the tutoring services was high, especially among those who used them. However, satisfaction was lower regarding the range of subjects offered and the availability of tutors.

\Large\textbf{Recommendations}

\normalsize

    \textbf{Communication Strategy:}
Consider enhancing digital communications through institutional channels like email, Moodle, and the CNC website. Maintaining a balance of digital and personal communications could help maximize reach and effectiveness.

    \textbf{Service Availability:}
        Explore extending tutoring hours, including evening and weekend sessions, and offering online tutoring to better accommodate students’ schedules. Increasing the cap on tutoring sessions per week might help meet higher demand during peak times. Additionally, expanding the range of subjects covered by tutoring services and developing online tutoring options and flexible scheduling could improve accessibility.

    \textbf{Awareness and Familiarity:}
        Orientation sessions and tutorial resources could guide students on how to access and use tutoring services. Promoting the benefits of tutoring, such as improving grades (promotion focus) rather than just avoiding failure (prevention focus), might also be beneficial.

    \textbf{Pilot Programs:}
         Testing the implementation of peer tutoring groups and additional tutoring sessions for high-demand subjects could provide valuable insights. Evaluating the impact of these pilot programs on student engagement and satisfaction might further inform service improvements.

\Large\textbf{Conclusion}

\normalsize

The survey reveals a clear demand for more accessible and flexible tutoring services at CNC. By addressing the identified barriers and enhancing communication strategies, the Academic Success Centre can better serve its student population. Considering the suggestions provided may increase the usage and satisfaction of tutoring services, fostering a more supportive and effective learning environment for CNC students.

\normalsize
  \newpage
}

\makeatother
\usepackage{scrlayer-scrpage}  %# For customizing page headers and footers.
\cohead[]{\vspace*{-.5cm}\includegraphics[width=7cm]{../fig/irem_logo.png}}  %# Logo in the header.
\pagenumbering{arabic}  %# Page numbering style.

\ifoot[]{For internal use only}  %# Footer information.
\cfoot[]{Institutional Research and Enrolment Management - CNC}  %# Footer information.
\ofoot[\pagemark]{\pagemark}  %# Page number formatting.
\setkomafont{pagefoot}{\footnotesize}  %# Font size for footer.
\usepackage[left=1.5cm,top=3cm,right=1.5cm,bottom=2.5cm]{geometry}        
\usepackage[absolute,overlay]{textpos}  %# For positioning images.
\usepackage{caption} %# For customizing captions.
\usepackage[labelfont=bf]{caption}


%% Use some custom fonts
\setsansfont{OpenSans}[
    Path=OpenSans/,
    Scale=0.9,
    Extension = .ttf,
    UprightFont=*-Regular,
    BoldFont=*-Bold,
    ItalicFont=*-Italic,
    ]
\setmainfont{OpenSans}[
    Path=OpenSans/,
    Scale=0.9,
    Extension = .ttf,
    UprightFont=*-Regular,
    BoldFont=*-Bold,
    ItalicFont=*-Italic,
    ]
    
\usepackage{etoolbox}
\AtBeginEnvironment{figure}{\addtocontents{lof}{\vspace{5pt}}}

\usepackage[none]{hyphenat}

